\documentstyle[%


righttag%


,11pt%


,amssymb%


,russian%               remove in Eng.


,izvru%                 Set izven% in Eng.


]{amsart}





% NB! \text{...} constructions





\newcommand{\la}{\langle}


\newcommand{\ra}{\rangle}


\renewcommand{\leq}{\leqslant}


\renewcommand{\le}{\leqslant}


\renewcommand{\ge}{\geqslant}


\renewcommand{\geq}{\geqslant}


\newcommand{\tts}[1]{}





\theoremstyle{plain}


\theorembodyfont{\it}


\newtheorem{alem}{\hskip\parindent Лемма}


%% \renewcommand{\thealem}{\Alph{alem}}  % alphbetic enumeration. eng


\renewcommand{\thealem}{\Asbuk{alem}}  % alphbetic enumeration russ.





\theoremstyle{definition}


\theorembodyfont{\rm}


\newtheorem{notenonum}{\hskip\parindent Замечание}


\renewcommand{\thenotenonum}{}





%\hoffset=-15mm%                  For Eng.


\setcounter{page}{13}


\god{2005}


\nomer{3~(514)} %                        Remove in Eng.


%\nomer{Vol.\,49, No.\,3}%               Set in Eng.


%\str{\pageref{begin}--\pageref{end}}%  Set in Eng.


\udk{517.518}





\author{\it Г.А.\,Акишев}


% NB:   ^^ remove in Eng.


\title{О порядках приближения функциональных классов


полиномами по обобщенной системе Хаара}


\thanks{}





\date{}


%\pagestyle{myheadings}%         Set in Eng.


\pagestyle{plain} %              Remove in Eng.


\begin{document}


%\thispagestyle{plain}%          Set in Eng.


\maketitle


\label{begin}








Пусть $\Bbb R^{d}$ --- $d$-мерное вещественное евклидово пространство, 


$\ov x=(x_{1},\dotsc,x_{d})$, $I^{d}=\{\ov x\in \Bbb 


R^{d}$; $0\leq x_{j}\leq 1$; $j=1,\dotsc,d\}$.





Банахово пространство $X$ измеримых по Лебегу на  $I^{d}$ функций 


называется симметричным, если


\begin{itemize}


\item[1.] из того, что $|f(\ov x)|\leq |g(\ov x)|$ почти всюду на $I^{d}$ 


и $g\in X$ следует, что $f\in X$ и $\|f\|_{X}=\|g\|_{X}$;


\item[2.] из $f\in X$ и равноизмеримости функций $|f(\ov x)|$ и $|g(\ov x)|$ 


следует $g\in X$ и $\|f\|_{X}=\|g\|_{X}$ (см. [1], с.\,123).


\end{itemize}


Здесь и в дальнейшем $\|f\|_{X}$ означает норму элемента $f\in X$.





Пусть $\chi_{e}(t)$ --- характеристическая функция множества $e\subset 


I^{d}$. Функция $\varphi(\mu e)=\|\chi_{e}\|_{X}$ называется 


фундаментальной функцией пространства $X$, где  $\mu e$ --- мера Лебега 


измеримого множества $e$. 


Таким образом, фундаментальная функция симметричного пространства $X$ есть 


функция $\varphi(t)=\|\chi_{[0,t]}\|_{X}$, определенная на отрезке $[0,1]$. 


Фундаментальную функцию  $\varphi(t)$ симметричного пространства можно 


считать вогнутой, неубывающей, непрерывной на $[0,1]$ функцией, причем 


$\varphi(0)=0$ ([1], с.\,137). Такие функции называются


$\Phi$-функциями.  Далее $X(\varphi)$ означает симметричное пространство с 


фундаментальной функцией $\varphi$. 





Ассоциированное к симметричному пространству $X(\varphi)$  пространство 


$X^{1}$ состоит из всех измеримых функций $g(t)$, для которых


$$


\|g\|_{X^{1}}=\sup\begin{Sb} f\in X \\ \|f\|_{X}\leq 1\end{Sb}


\int_{0}^{1}f(t)g(t)dt.


$$





Известно, что сепарабельность симметричного пространства $X(\varphi)$ 


является необходимым и достаточным условием совпадения ассоциированного к 


нему пространства $X^{1}$ со всем сопряженным пространством 


$X'(\ov\varphi)$ ([1], с.\,138), при этом 


$\ov{\varphi}(t)=\frac{t}{\varphi(t)}$, $t\in (0,1]$ и 


$\ov{\varphi}(0)=0$. 





Далее будем рассматривать сепарабельные симметричные пространства.


Сепарабельными симметричными пространствами являются, например,





1. $L_{q}(I^{d})$ --- пространство Лебега с нормой 


$$


\|f\|_{q}=\bigg(\int_{I^{d}}|f(\ov x)|^{q}d\ov 


x\bigg)^{\frac{1}{q}},\quad 1<q<+\infty;


$$





2. пространство Марцинкевича $M_{q}(I^{d})$ с нормой 


$$


\|f\|_{q,\infty}=\sup_{t\in(0,1]}t^{\frac{1}{q}-1}\int_{0}^{t}f^{*}


(\tau)d\tau,


$$


где $f^{*}(\tau)$ --- невозрастающая перестановка функций $|f(\ov x)|$ (см. 


[1], с.\,83);





3. пространство Лоренца $L_{q\theta}(I^{d})$ с нормой 


$$


\|f\|_{q\theta}=\bigg\{\frac{\theta}{q}\int_{0}^{1}\bigg(\int_{0}^{t}f^{*}


(\tau)d\tau\bigg)^{\theta}t^{\theta(\frac{1}{q}-1)-1}dt


\bigg\}^{\frac{1}{\theta}}<+\infty, \quad 1\leq q<+\infty, \quad 


1<\theta<+\infty.


$$





Для функции $\Psi(t)$, $t\in [0,1]$, положим


$$


\alpha_{\Psi}=\underset{t\rightarrow 0}{\varliminf}{}


\frac{\Psi(2t)}{\Psi(t)},\quad 


\beta_{\Psi}=\underset{t\rightarrow 0}{\varlimsup}{}\frac{\Psi(2t)}{\Psi(t)}.


$$


Известно, что для любого симметричного пространства $X(\varphi)$ справедливы 


неравенства $1\leq\alpha_{\varphi}\leq\beta_{\varphi}\leq 2$.





Пусть


$$


\Omega(f,\ov t)_{X}\equiv\Omega(f,t_{1},\dotsc,t_{d})_{X} 


=\sup\begin{Sb} 0\leq h_{j}\leq t_j \\ j=1,\dotsc,d \end{Sb}


\bigg\|\Delta_{\ov 


h}^{\ov 1}f\prod_{j=1}^{d}\chi_{[0,1-h_{j}]}(x_{j})\bigg\|_{X}


$$


--- смешанный модуль непрерывности функции $f\in X(\varphi)$, где 


$\Delta_{\ov h}^{\ov 1}f(\ov 


x)=\Delta_{h_{1}}^{1}(\dotsc\Delta_{h_{d}}^{1}f(x_{1},\dotsc,x_{d}))$


--- смешанная разность функции $f\in X(\varphi)$. 





Рассмотрим функциональный класс 


$$


H_{X}^{\ov r}=\bigg\{f\in X(\varphi):\Omega(f,\ov 


t)_{X(\varphi)}\leq\prod_{j=1}^{d}t_{j}^{r_{j}},\ \ov{t}\in I^{d}\bigg\}, 


\ \ 0<r_{j}\leq 1\ \; \text{для всех}\ \; j=1,\dotsc,d.


$$


В случае $r_{1}=\dotsb=r_{d}=r$ вместо 


$H_{X}^{\ov r}$ будем писать $H_{X}^{r}$.





Через $C(q,p,r,\dotsc)$ обозначим величины, 


зависящие от указанных параметров.


Запись $A\asymp B$ означает, что существуют положительные постоянные $c_{1}$, 


$c_{2}$ такие, что $c_{1}A\leq B\leq c_{2}A$.





Рассмотрим обобщенную систему Хаара, определенную на отрезке $[0,1]$. Пусть 


дана последовательность $\{p_{n}\}$ натуральных чисел, $p_{n}\geq 2$ для всех 


$n=1,2,\dotsc$ Обобщенная система Хаара $\chi\{p_{n}\}=\{\chi_{n}(t)\}$ 


определяется следующим образом ([2], [3]): положим 


$\chi_{1}(t)\equiv 1$ на отрезке $[0,1]$. Заданное натуральное число $n\geq 


2$ представляется в виде $n=m_{k}+r(p_{k+1}-1)+s$, где $m_{k}=p_{1}\cdot 


p_{2}\cdot\dotsc\cdot p_{k}$; $k=1,2,\dotsc$; $r=0,1,\dotsc,m_{k}-1$;


$s=1,2,\dotsc,p_{k+1}-1$.





Через $A$ обозначим множество точек вида $\frac{l}{m_{k}}$ на отрезке 


$[0,1]$, $l=0,1,\dotsc$ Тогда при $t\in B\equiv [0,1]\setminus A$ разложение 


$t=\sum\limits_{k=1}^{\infty}\frac{\alpha_{k}(t)}{m_{k}}$,


$\alpha_{k}(t)=0,1,\dotsc,p_{k}-1$, единственно. Далее, определим функцию 


$\chi_{n}(t)\equiv\chi_{k,r}^{s}(t)$ следующим образом:


$$


\chi_{n}(t)=\chi_{k,r}^{s}(t)=


\begin{cases}


\sqrt{m_{k}} \exp\big\{\frac{2\pi 


is\alpha_{k+1}(t)}{p_{k+1}}\big\}, & t\in 


\big(\frac{r}{m_{k}}, \frac{r+1}{m_{k}}\big)\cap B;\\[2mm]


0, & t\notin \big[\frac{r}{m_{k}}, 


\frac{r+1}{m_{k}}\big].


\end{cases}


$$


Пользуясь тем, что множество $B$ всюду плотно на $[0,1]$, функцию 


$\chi_{n}(t)$  по непрерывности продолжим на интервал 


$\big(\frac{r}{m_{k}}, \frac{r+1}{m_{k}}\big)$. Затем в точках разрыва 


эту функцию положим равной полусумме ее предельных значений справа и слева, 


а на концах 0 и 1 --- ее предельным значениям изнутри отрезка.


Определенная таким образом система $\chi\{p_{n}\}$ ортонормирована и полна в 


пространстве $L_{1}$ ([2], [3]).





Пусть даны $\{p_{n_{j}}^{(j)}\}$ --- последовательности натуральных чисел


$p_{n_{j}}^{(j)} \geq 2$; $n_{j}=1,2,\dotsc$;\linebreak $j=1,2,\dotsc,d$ и 


$m_{n_{j}}^{(j)}=p_{1}^{(j)}\cdot\dotsc\cdot p_{n_{j}}^{(j)}$. Через 


$\{\chi_{\ov n}(\ov 


x)\}=\Big\{\prod\limits_{j=1}^{d}\chi_{n_{j}}(x_{j})\Big\}$ обозначим 


кратную обобщенную систему Хаара, $a_{\ov n}(f)$ --- коэффициенты Фурье 


функции $f\in L_{1}(I^{d})$ по этой системе.


Пусть также


$$


T_{\ov n}(\ov 


x)=\sum_{k_{1}=1}^{n_{1}}\dotsb\sum_{k_{d}=1}^{n_{d}}a_{k_{1},\dotsc,k_{d}} 


\prod_{j=1}^{d}\chi_{k_{j}}(x_{j})=\sum_{\ov 1\leq\ov k\leq\ov n}a_{\ov 


k}\chi_{\ov k}(\ov x)


$$


--- полином по обобщенной системе Хаара порядка $n_{j}$ по переменной $x_{j}$. 


Здесь неравенство $\ov k\leq\ov n$ понимается в том смысле, что 


$k_{j}\leq n_{j}$ для всех $j=1,\dotsc,d$. 





Положим 


\begin{gather*}


\delta_{\ov \nu}(f,\ov x)=\sum_{n_{1}= 


m_{\nu_{1}}^{(1)}}^{m_{\nu_{1}+1}^{(1)}-1}\dotsb\sum_{n_{d}= 


m_{\nu_{d}}^{(d)}}^{m_{\nu_{d}+1}^{(1)}-1}a_{n_{1},\dotsc,n_{d}}(f) 


\prod_{j=1}^{d}\chi_{k_{j}}(x_{j}),\\


%


S_{n}^{\ov\gamma}(f,\ov{x})=\sum_{\la\ov{s},\ov{\gamma}\ra\leq 


n}\delta_{\ov s}(f,\ov{x}),


\end{gather*}


где $\la\ov s,\ov{\gamma}\ra=\sum\limits_{j=1}^{d}s_{j}\gamma_{j}$.  В случае 


$\ov{\gamma}=(1,\dotsc,1)$ будем писать $S_{n}(f,\ov{x})$ и $\|\ov s\|$ 


соответственно вместо $S_{n}^{\ov\gamma}(f,\ov{x})$ и 


$\la\ov{s},\ov{\gamma}\ra$.





Порядки приближения классов $H_{\theta}^{r}$ по норме пространства Лебега 


$L_{q}(I^{d})$ полиномами по кратной системе Хаара с гармониками из 


гиперболических крестов исследованы в [4], а оценки $N$-членных 


приближений --- в [5]--[7]. 





В данной статье изучаются порядки приближения функциональных классов 


$H_{X}^{r}$ полиномами по обобщенной системе Хаара по норме 


пространства Марцинкевича.








\begin{alem}[{\series{m}\shape{n}\selectfont [8], лемма 4}]


Пусть даны $\Phi$-функции 


$\phi_{1}(x)$ и $\phi_{2}(x)$, $x\in (0,1]$, и 


$\beta_{\phi_{1}}<\alpha_{\phi_{2}}$. Тогда для функции


$$


\theta(x)=


\begin{cases}


\frac{\phi_{2}(x)}{\phi_{1}(x)}, & x\in(0,1]; \\


0, & x=0,


\end{cases}


$$


существует $\Phi$-функция $\theta_{1}(x)$, для которой 


$\alpha_{\theta_{1}}>1$ и $\theta_{1}(x)\asymp\theta(x)$, $x\in[0,1]$.


\end{alem}








\begin{alem}


Для любых чисел $p>1$, $\alpha>0$ имеет место соотношение


$$


\sum_{\la\ov{s},\ov{\gamma}\ra >n}p^{-\alpha\la\ov{s},\ov{\gamma}\ra}\asymp 


p^{-\alpha n}\cdot n^{d-1}.


$$


\end{alem}








\begin{alem}


Пусть $\ov{\gamma}=(\gamma_{1},\dotsc,\gamma_{d})$ такие, 


что $1=\gamma_{j}$, $j=1,\dotsc,l;$ $1<\gamma_{j}$, $j=l+1,\dotsc,d$.


Тогда для любых чисел $p>1$, $\alpha>0$ имеет место неравенство


$$


\sum_{\la\ov{s},\ov{\gamma}\ra\leq n}p^{\alpha\la\ov{s},\ov{\gamma}\ra}\leq 


C(\alpha,p)\cdot p^{\alpha n}\cdot n^{d-1}.


$$


\end{alem}





Леммы Б, В доказываются так же, как в случае $p=2$ ([9], с.\,11). 





Теперь докажем несколько вспомогательных утверждений.  


   





\begin{lem}


Для любой функции  $f\in X(\varphi)$ имеет место неравенство


$$


\|\delta_{\ov\nu}(f)\|_{X}\leq 


C(d)\big\|\Delta_{m^{-1}_{\nu_1+1},\dotsc,m^{-1}_{\nu_d+1}} ^{\ov 


1}f\big\|_{X}\cdot\prod_{j=1}^{d}p_{\nu_{j}+1}^{(j)}


\ln p_{\nu_{j}+1}^{(j)}.


$$


\end{lem}








{\bf Доказательство} проведем для случая $d=2$. Учитывая определение 


функций обобщенной системы Хаара и применяя преобразование Абеля, получим


\begin{multline}


a_{n_{1}n_{2}}(f)=\sqrt{m_{\nu_{1}}^{(1)}\cdot m_{\nu_{2}}^{(2)}}


\sum_{k_{1}=1}^{p_{\nu_{1}+1}^{(1)}-1}\sum_{k_{2}=1}^{p_{\nu_{2}+1} 


^{(2)}-1}\int_{\rho_{\nu_{1}r_{1}k_{1}}}\int_{\rho_{\nu_{2}r_{2}k_{2}}} 


\Delta_{m_{\nu_{1}+1}^{-1}m_{\nu_{2}+1}^{-1}}^{11}f(x_{1},x_{2})dx_{1}dx_{2}


\times \\


%


\times\frac{1-\exp\{-2\pi i\frac{k_{1}}{p_{\nu_{1}+1}^{(1)}}\}}


{1-\exp\{-2\pi i\frac{s_{1}}{p_{\nu_{1}+1}^{(1)}}\}}\cdot 


\frac{1-\exp\{-2\pi i\frac{k_{2}}{p_{\nu_{2}+1}^{(2)}}\}}


{1-\exp\{-2\pi i\frac{s_{2}}{p_{\nu_{2}+1}^{(2)}}\}},


\end{multline}


где


$$


\rho_{\nu rk}=\bigg(\frac{r}{m_{\nu}}+\frac{k-1}{m_{\nu+1}},


\frac{r}{m_{\nu}}+\frac{k}{m_{\nu+1}}\bigg).


$$





Положим


$$


I(\ov r)=\bigg[\frac{r_{1}}{m_{\nu_{1}}},\frac{r_{1}+1}{m_{\nu_{1}}} \bigg]


\times\bigg[\frac{r_{2}}{m_{\nu_{2}}},\frac{r_{2}+1}{m_{\nu_{2}}} \bigg].


$$





В силу равенства (1), определения функций $\chi_{\nu,r}^{(s)}$ и свойства 


интеграла для любой функции $g\in X'(\varphi)$, $\|g\|_{X'}\leq 1$, имеем


\begin{multline}


\bigg|\int_{0}^{1} \int_{0}^{1}\delta_{\ov \nu}(f,\ov t)g(\ov 


t)dt_{1}dt_{2}\bigg|=


\bigg|\int_{0}^{1} \int_{0}^{1}\sum_{r_{1}=0}^{m_{\nu_{1}}^{(1)}-1}


\sum_{r_{2}=0}^{m_{\nu_{2}}^{(2)}-1} \sum_{s_{1}=1}^{p_{\nu_{1}+1}^{(1)}-1} 


\sum_{s_{2}=1}^{p_{\nu_{2}+1}^{(2)}-1}a_{\ov\nu,\ov r}^{(\ov s)}(f) 


\prod_{j=1}^{2}\chi_{\nu_{j},r_{j}}^{(s_{j})}(t_{j})g(\ov 


t)dt_{1}dt_{2}\bigg|\leq \\


%


\leq 4m_{\nu_{1}}^{(1)}\cdot m_{\nu_{2}}^{(2)}


\int_{0}^{1} \int_{0}^{1}\sum_{r_{1}=0}^{m_{\nu_{1}}^{(1)}-1}


\sum_{r_{2}=0}^{m_{\nu_{2}}^{(2)}-1} \sum_{s_{1}=1}^{p_{\nu_{1}+1}^{(1)}-1} 


\sum_{s_{2}=1}^{p_{\nu_{2}+1}^{(2)}-1} 


\prod_{j=1}^{2}\bigg|1-\exp\big\{-2\pi 


i\frac{s_{j}}{p_{\nu_{j}+1}^{(j)}}\big\}\bigg|^{-1}\times \\


%


\times\bigg[\sum_{k_{1}=1}^{p_{\nu_{1}+1}^{(1)}-1} 


\sum_{k_{2}=1}^{p_{\nu_{2}+1}^{(2)}-1} 


\int_{\rho_{\nu_{1}r_{1}k_{1}}}\int_{\rho_{\nu_{2}r_{2}k_{2}}} 


\big|\Delta_{(m_{\nu_{1}+1}^{(1)})^{-1}(m_{\nu_{2}+1}^{(2)})^{-1}}^{11}


f(x_{1},x_{2})\big|dx_{1}dx_{2} \bigg]\chi_{I(\ov r)}(\ov t)g(\ov t)dt_{1} 


dt_{2}.


\end{multline}


Далее, учитывая неравенство ([3], с.\,307)


$$


\sum_{s=1}^{p_{\nu+1}-1}\bigg|1-\exp\bigg\{-2\pi 


i\frac{s}{p_{\nu+1}}\bigg\}\bigg|^{-1}\leq C\cdot p_{\nu+1} \ln p_{\nu+1},


$$


из (2) получим


\begin{multline*}


\bigg|\int_{0}^{1} \int_{0}^{1}\delta_{\ov \nu}(f,\ov t)g(\ov 


t)dt_{1}dt_{2}\bigg|\leq


\sum_{r_{1}=0}^{m_{\nu_{1}}^{(1)}-1}


\sum_{r_{2}=0}^{m_{\nu_{2}}^{(2)}-1}\int_{I(\ov r)}\bigg[\sup_{\ov t\in 


I} \frac{1}{|I|}\int_{I}\big| 


\Delta_{(m_{\nu_{1}+1}^{(1)})^{-1}(m_{\nu_{2}+1}^{(1)})^{-1}}^{11} 


f(x_{1},x_{2})\big|\times\\


%


\times\chi_{\rho_{\nu_{1}r_{1}}^{-}\rho_{\nu_{2}r_{2}}^{-}}(x_{1},x_{2}) 


dx_{1}dx_{2}\bigg]\cdot |g(\ov t)|dt_{1}dt_{2}=


4\prod_{j=1}^{2}p_{\nu_{j}+1}^{(j)} \ln p_{\nu_{j}+1}^{(j)}\int_{0}^{1} 


\int_{0}^{1}\bigg[\sup_{\ov t\in 


I} \frac{1}{|I|}\times \\


%


\times\int_{I}\big| 


\Delta_{(m_{\nu_{1}+1}^{(1)})^{-1}(m_{\nu_{2}+1}^{(1)})^{-1}}^{11} 


f(x_{1},x_{2})\big|\chi_{\rho_{\nu_{1}r_{1}}^{-}\rho_{\nu_{2}r_{2}}^{-}}


(x_{1},x_{2}) dx_{1}dx_{2}\bigg]\cdot |g(\ov t)|dt_{1}dt_{2}


\end{multline*}


(где $\rho_{\nu r}= \big[\frac{r}{m_{\nu}}, 


\frac{r+1}{m_{\nu}}-\frac{1}{m_{\nu+1}}\big]$) для любой функции 


$g\in X'(\ov \varphi)$, $\|g\|_{X'}\leq 1$.





Следовательно, в силу ограниченности максимального оператора Харди в 


симметричном пространстве ([1], с.\,187) будем иметь


$$


\|\delta_{\ov\nu}(f)\|_{X}\leq 


4\prod_{j=1}^{2}p_{\nu_{j}+1}^{(j)}\ln p_{\nu_{j}+1}^{(j)}\cdot


\big\|\Delta_{m_{\nu_1+1}^{-1},m_{\nu_2+1}^{-1}}^{11}


f\cdot\chi_{\rho} \big\|_{X}.\quad\square


$$








\begin{notenonum}


Из доказанной леммы в случае $X(\varphi)=L_{q}(I^{d})$, $1\leq q<+\infty$, 


и $p_{n_{j}}^{(j)}=2$ \, $\forall n_{j}=1,2,\dotsc$; $j=1,\dotsc,d$ (т.\,е. 


$\chi_{n}\{p_{n}\}$ --- система Хаара), следует лемма 1 [4].


\end{notenonum}








\begin{lem}


Пусть $X(\varphi)$ --- симметричное пространство, 


$1<\alpha_{\varphi},\beta_{\varphi}\leq 2$. Тогда для любого полинома 


$T_{{\ov m}_{n}}(\ov x)$ по обобщенной системе Хаара имеет место 


неравенство 


$$


\|T_{{\ov m}_{n}}\|_{\infty}\leq C(\varphi)\cdot 


\frac{1}{\varphi\Big[\Big(\prod\limits_{j=1}^{d}m_{n_{j}}^{(j)}\Big)^{-1} 


\Big]}\cdot \|T_{{\ov m}_{n}}\|_{X}.


$$


\end{lem}








{\bf Доказательство.} Известно, что ([10])


$$


T_{{\ov m}_{n}}(\ov y)=\int_{I^{d}}T_{{\ov m}_{n}}(\ov 


x)\cdot\prod_{j=1}^{d}D_{m_{n_{j}}^{(j)}}(x_{j},y_{j})d\ov x,


$$


где $D_{n}(u,t)=\sum\limits_{l=1}^{n}\chi_{l}(u){\ov\chi}_{l}(t)$ --- ядро 


Дирихле по обобщенной системе Хаара, и ([3], с.\,300)


$$


D_{m_{n}}(u,t)=


\begin{cases}


m_{n}, & \text{если} \ \ u,t\in\big(\frac{r}{m_{n}}, 


\frac{r+1}{m_{n}}\big);\\


0 & \text{в остальных случаях}.


\end{cases}


$$


Из этих равенств и определения максимального пространства следует 


утверждение леммы.





Отметим, что лемма 2 в одномерном случае доказана в [11], а в случае   


$X(\varphi)=L_{p}(I^{d})$ (пространство Лебега), $1\leq p<+\infty$,


$d\geq 1$, доказана в [10].








\begin{lem}


Пусть $X(\varphi)$ --- симметричное пространство,  


$1<\alpha_{\varphi},\beta_{\varphi}\leq 2$, и кратная обобщенная система 


Хаара определена ограниченными последовательностями $\{p_{n}^{(j)}\},


p_{n}^{(j)}\leq p^{(j)}$, $j=1,\dotsc,d$,


$p=\max\limits_{j}p^{(j)}$. Тогда для 


любой функции  $f\in X(\varphi)$ имеет место неравенство


$$


\frac{1}{x}\int_{0}^{x}f^{*}(t)dt\leq 


C(\varphi,p)\bigg\{\frac{1}{\varphi(x)}\cdot 


\|f-S_{n}^{(\ov\gamma)}(f)\|_{X}+\sum_{\la\ov{s},\ov{\gamma}\ra\leq 


n}\bigg[\varphi\bigg(\bigg[\prod_{j=1}^{d}m_{s_{j}}^{(j)}\bigg]^{-1} 


\bigg)\bigg]^{-1}\|\delta_{\ov s}(f)\|_{X}\bigg\}


$$


для каждого $x\in (p^{-n},p^{-n+1}]$, $n=1,2,\dotsc$ 


\end{lem}





\begin{pf}


Воспользуемся известным равенством ([1], с.\,89)


$$


\frac{1}{x}\int_{0}^{x}f^{*}(t)dt=\frac{1}{x}


\sup\begin{Sb} A\subset I^d \\ \mu A=x \end{Sb}


\int_{A}|f(\ov x)|d\ov x, \quad x\in (0,1].


$$


Тогда 


\begin{equation}


\frac{1}{x}\int_{0}^{x}f^{*}(t)dt\leq \frac{1}{x}\bigg\{


\sup\begin{Sb} A\subset I^d \\ \mu A=x \end{Sb}


\int_{A}|f(\ov x)-S_{n}^{(\ov\gamma)}(f,\ov x)|d\ov x+


\sup\begin{Sb} A\subset I^d \\ \mu A=x \end{Sb}


\int_{A}|S_{n}^{(\ov\gamma)}(f,\ov x)| d\ov x\bigg\}.


\end{equation}


Применяя неравенство Г\"ельдера (см. определение максимального симметричного 


пространства), получим


\begin{equation}


\int_{A}|f(\ov x)-S_{n}^{(\ov\gamma)}(f,\ov x)|d\ov x\leq \frac{\mu A}{{\ov 


\phi}(\mu A)}\cdot \|f-S_{n}^{(\ov\gamma)}(f)\|_{X}


\end{equation}


для любого измеримого множества $A\subset I^{d}$. Далее, пользуясь 


неравенством разных метрик (см. лемму 2), будем иметь


\begin{multline}


\int_{A}|S_{n}^{(\ov\gamma)}(f,\ov x)| d\ov 


x\leq\sum_{l=1}^{n}\;\sum_{l-1<\la\ov{s},\ov{\gamma}\ra\leq l}


\int_{A}|\delta_{\ov s}(f,\ov x)| d\ov x \leq \\


%


\leq C(\varphi,d) 


\sum_{l=1}^{n}\;\sum_{l-1<\la\ov{s},\ov{\gamma}\ra\leq l}\mu 


A\bigg[\varphi\bigg(\bigg[ 


\prod_{j=1}^{d}m_{s_{j}}^{(j)}\bigg]^{-1}\bigg)\bigg]^{-1}


\|\delta_{\ov s}(f)\|_{X}.


\end{multline}


Из неравенств (3)--(5) следует утверждение леммы.


\end{pf}








\begin{lem}


Пусть $0<r<\log_{2}\alpha_{\varphi}<1$. Тогда функция одной переменной


$$


f_{0}(t)=\sum_{k=1}^{\infty}m_{k}^{-(r+\frac{1}{2})}\cdot 


\frac{1}{\varphi(m_{k}^{-1})}\chi_{m_{k}+p_{k+1}}(t), \quad t\in[0,1],


$$


принадлежит классу $H_{X}^{r}$.


\end{lem}








\begin{pf}


Так как 


$\|\chi_{m_{k}+p_{k+1}}\|_{X}=\sqrt{m_{k}}\cdot 


\varphi(m_{k}^{-1})$, то


$$


\|f_{0}\|_{X}\leq\sum_{k=1}^{\infty} \frac{m_{k}^{-(r+\frac{1}{2})}}{ 


\varphi(m_{k}^{-1})}\|\chi_{m_{k}+p_{k+1}}\|_{X}=\sum_{k=1}^ 


{\infty}m_{k}^{-r}<+\infty,


$$


т.\,е. $f_{0}\in X(\varphi)$.





Докажем, что $C_{0}\cdot f_{0}\in H_{X}^{r}$, где $C_{0}$ 


--- некоторое положительное число. Пусть $h\in (0,1]$ и 


натуральное число $\nu$ такое, что $h\in (m_{\nu+1}^{-1}, m_{\nu}^{-1}]$. 


Тогда по свойству нормы имеем


\begin{multline}


\|[f(\cdot+h)-f(\cdot)]\chi_{[0, 1-h]}\|_{X}\leq


\sum_{k=1}^{\nu} \frac{m_{k}^{-(r+\frac{1}{2})}}{ 


\varphi(m_{k}^{-1})}\|[\chi_{m_{k}+p_{k+1}}(\cdot+h)-\chi_{m_{k}+p_{k+1}}]  


\chi_{[0,1-h]} \|_{X}+ \\


%


+2\sum_{k=\nu+1}^{\infty} \frac{m_{k}^{-(r+\frac{1}{2})}}{ 


\varphi(m_{k}^{-1})}\|\chi_{m_{k}+p_{k+1}}\|_{X}.


\end{multline}


Так как $h\in (m_{\nu+1}^{-1}, m_{\nu}^{-1}]$, то $1-h\geq 


1-\frac{1}{m_{\nu}}>1-\frac{1}{p_{1}}>\frac{2}{p_{1}}$ при $p_{1}\geq 3$. 


Следовательно, по определению функции $\chi_{n}(t)$ имеем  


$\chi_{m_{k}+p_{k+1}}(t)=\chi_{m_{k}+p_{k+1}}(t+h)=0$ для 


$t\in(\frac{2}{p_{1}}, 1-h]$. Поэтому для чисел $k=1,2,\dotsc,\nu$ имеют место 


равенства 


\begin{multline}


\int\limits_{0}^{1-h}[\chi_{m_{k}+p_{k+1}}(t+h)-\chi_{m_{k}+p_{k+1}}(t)] 


g(t)dt=\int\limits_{0}^{2/p_1}[\chi_{m_{k}+p_{k+1}}(t+h)-


\chi_{m_{k}+p_{k+1}}(t)] g(t)dt=\\


%


=\sum_{j=0}^{p_{k+1}-1}\int\limits_{\frac{1}{m_{k}}+\frac{j}{m_{k+1}}}^{


\frac{1}{m_{k}}+\frac{j+1}{m_{k+1}}}


[\chi_{m_{k}+p_{k+1}}(t+h)-\chi_{m_{k}+p_{k+1}}(t)] g(t)dt.


\end{multline}


По выбору $h\leq\frac{1}{m_{\nu}}\leq\frac{1}{m_{k+1}}$, 


$k=1,\dotsc,\nu-1$. Поэтому если $t\in 


\big(\frac{1}{m_{k}}+\frac{j+1}{m_{k+1}}-h, \frac{1}{m_{k}}+ 


\frac{j+1}{m_{k+1}}\big)$, то $t+h\in 


\big(\frac{1}{m_{k}}+\frac{j+1}{m_{k+1}}, \frac{1}{m_{k}}+ 


\frac{j+2}{m_{k+1}}\big)$. Следовательно, по определению функций 


$\chi_{m_{k}+p_{k+1}}(t)$ имеем 


\begin{multline}


\int\limits_{\frac{1}{m_{k}}+\frac{j}{m_{k+1}}}^{


\frac{1}{m_{k}}+\frac{j+1}{m_{k+1}}}


[\chi_{m_{k}+p_{k+1}}(t+h)-\chi_{m_{k}+p_{k+1}}(t)] g(t)dt= \\


%


=\sqrt{m_{k}}  


\int\limits_{\frac{1}{m_{k}}+\frac{j+1}{m_{k+1}}-h}^{\frac{1}{m_{k}}+\frac{j+1}{m_{k+1}}}


\bigg[\exp\bigg\{\frac{2\pi i(j+1)}{p_{k+1}}\bigg\}-


\exp\bigg\{\frac{2\pi i}{p_{k+1}}\bigg\}\bigg]g(t)dt


\end{multline}


для любой функции $g\in X'(\ov\varphi)$, $\|g\|_{X'}\leq 1$. Из (7), 


(8) следует


\begin{equation}


\sum_{k=1}^{\nu} \frac{m_{k}^{-(r+\frac{1}{2})}}{ 


\varphi(m_{k}^{-1})}\|[\chi_{m_{k}+p_{k+1}}(\cdot+h)-\chi_{m_{k}+p_{k+1}}]  


\chi_{[0,1-h]} \|_{X}\leq C\cdot \varphi(h)\cdot\sum_{k=1}^{\nu}


\frac{m_{k}^{-r}}{\varphi(m_{k}^{-1})}. 


\end{equation}


По условию леммы $2^{r}<\alpha_{\varphi}$. Тогда $\exists 


\varepsilon>0:2^{r}<2^{r+\varepsilon}<\alpha_{\varphi}$. Поэтому в 


силу леммы~A для функции 


$$


\theta(t)=


\begin{cases}


\frac{\varphi(t)}{t^{r+\varepsilon}}, & t\in(0,1]; \\ 


0, & t=0,


\end{cases}


$$


существует $\Phi$-функция $\theta_{1}(t)$ такая, что $\theta(t)\asymp 


\theta_{1}(t)$. Значит, функция $\theta_{1}(t)=\frac{\varphi(t)}{t^{r+ 


\varepsilon}}$ возрастает на $(0,1]$. Поэтому


$$


\sum_{k=1}^{\nu}


\frac{m_{k}^{-r}}{\varphi(m_{k}^{-1})}\leq C\cdot 


\frac{m_{\nu}^{-r}}{\varphi(m_{\nu}^{-1})}.


$$


Тогда, учитывая $h\in (m_{\nu+1}^{-1}, m_{\nu}^{-1}]$ и свойства 


функции $\varphi$, из неравенства (9) получим


\begin{equation}


\sum_{k=1}^{\nu} \frac{m_{k}^{-(r+\frac{1}{2})}}{ 


\varphi(m_{k}^{-1})}\|[\chi_{m_{k}+p_{k+1}}(\cdot+h)-\chi_{m_{k}+p_{k+1}}]  


\chi_{[0,1-h]} \|_{X}\leq C\varphi(h)\cdot  


\frac{m_{\nu}^{-r}}{\varphi(m_{\nu}^{-1})}\leq C\cdot m_{\nu}^{-r}\leq 


C\cdot h^{r}.


\end{equation}


Далее, т.\,к. ([10]) $\|\chi_{n}\|_{X}\asymp 


\varphi(m_{k}^{-1})m_{k}^{\frac{1}{2}}$ для  $n=m_{k}+1,\dotsc,m_{k+1}$,


то 


\begin{equation}


\sum_{k=\nu+1}^{\infty} \frac{m_{k}^{-(r+\frac{1}{2})}}{ 


\varphi(m_{k}^{-1})}\|\chi_{m_{k}+p_{k+1}}\|_{X}  


\leq C\cdot \sum_{k=\nu+1}^{\infty} m_{k}^{-r}\leq C\cdot m_{\nu+1}^{-r}


C\cdot h^{r}. 


\end{equation}


Из неравенств (6), (10), (11) следует 


$$


\|[f_{0}(\cdot+h)-f_{0}(\cdot)]\chi_{[0, 1-h]}\|_{X}\leq C\cdot h^{r}.


$$


Таким образом, $\frac{1}{C}f_{0}\in H_{X}^{r}$.


\end{pf}





Далее положим $\gamma_{j}=\frac{r_{j}}{r_{1}}$, $j=1,\dotsc,d$.








\begin{thm}


Пусть обобщенная система Хаара определена ограниченными последовательностями 


$\{p_{n}^{(j)}\}$, $p_{n}^{j}\leq p^{(j)}$, $j=1,\dotsc,d$,


$p=\overline{\max}p^{(j)};$ $X(\varphi)$ --- симметричное пространство с 


фундаментальной функцией $\varphi$ и $\frac{1}{q}<\log_{2} 


\alpha_{\varphi}$, $\log_{2} \beta_{\varphi}<\frac{1}{q}+r_{1}$, $1\leq 


q<+\infty$, $0<r_{1}<1$, $0<r_{1}=\dotsb=r_{l}<r_{l+1}\leq\dotsb\leq r_{d}$. 


Если


\begin{equation}


\sup_{n>0}\bigg\{\frac{p^{-\frac{n}{q}}}{\varphi(p^{-n})}\sum_{\la\ov{s},


\ov{\gamma}\ra>n}\|\delta_{\ov s}(f)\|_{X}+p^{-\frac{n}{q}}


\sum_{\la\ov{s}, \ov{\gamma}\ra\leq n}\frac{1}{\varphi(p^{-\|\ov s\|})}


\|\delta_{\ov s}(f)\|_{X}\bigg\}<+\infty,


\end{equation}


то $f\in M_{q}(I^{d})$ и имеет место оценка


$$


\sup_{f\in H_{X}^{\ov r}}\|f-S_{n}^{(\ov\gamma)}(f)\|_{q,\infty}


\leq C\cdot \frac{p^{-n(r_{1}+\frac{1}{q})}}{\varphi(p^{-n})}\cdot n^{d-1}.


$$


\end{thm}








{\bf Доказательство.} Пользуясь леммой 3, для любого $x\in (p^{-n}, 


p^{-n+1}]$ получим


$$


x^{\frac{1}{q}-1}\int_{0}^{x}f^{*}(t)dt\leq C(p,d,\varphi) 


\sup_{n>0}\bigg\{\frac{p^{-\frac{n}{q}}}{\varphi\big(\frac{1}{p^{n}}\big)} 


\|f-S_{n}^{(\ov\gamma)}(f)\|_{X}+


p^{-\frac{n}{q}}\sum_{0<\la\ov{s},\ov{\gamma}\ra\leq n} 


\frac{1}{\varphi(p^{-\|\ov s\|})}\|\delta_{\ov s}(f)\|_{X}\bigg\}.


$$


Следовательно, в силу условия (12) $f\in M_{q}(I^{d})$ и 


\begin{equation}


\|f\|_{q,\infty} \leq C(p,d,\varphi) 


\sup_{n>0}\bigg\{\frac{p^{-\frac{n}{q}}}{\varphi(p^{-n})}\sum_{\la\ov{s},


\ov{\gamma}\ra>n}\|\delta_{\ov s}(f)\|_{X}+ p^{-\frac{n}{q}}


\sum_{0<\la\ov{s},\ov{\gamma}\ra\leq n} 


\frac{1}{\varphi(p^{-\|\ov s\|})}\|\delta_{\ov s}(f)\|_{X}\bigg\}.


\end{equation}


Применяя неравенство (13) к функции $f-S_{n}^{(\ov\gamma)}(f)


\in M_{q}(I^{d})$, будем иметь


\begin{multline}


\|f-S_{n}^{(\ov\gamma)}(f)\|_{q,\infty}\leq C(p,d,\varphi) 


\sup_{\nu>0}\bigg\{\frac{p^{-\frac{\nu}{q}}}{\varphi(p^{-\nu})}\sum_{\la\ov{s}, 


\ov{\gamma}\ra>\nu}\|\delta_{\ov s}(f-S_{n}^{(\ov\gamma)}(f))\|_{X}+ \\


%


+p^{-\frac{1}{q}}\sum_{\la\ov{s},\ov{\gamma}\ra\leq \nu} 


\frac{1}{\varphi(p^{-\|\ov s\|})}\|\delta_{\ov 


s}(f-S_{n}^{(\ov\gamma)}(f))\|_{X}\bigg\}.


\end{multline}


Если


$\la\ov{s},\ov{\gamma}\ra\leq n$, то $ \delta_{\ov 


s}(f-S_{n}^{(\ov\gamma)}(f))(x)=0$. Если же 


$\la\ov{s},\ov{\gamma}\ra> n$, то $ \delta_{\ov 


s}(f-S_{n}^{\ov\gamma}(f))(\ov x)=\delta_{\ov s}(f,\ov x)$.


Учитывая эти равенства, получим 


\begin{gather}


\sum_{\la\ov{s},\ov{\gamma\ra}\le\nu}\frac{1}{\varphi(p^{-\|\ov{s}\|})} 


\|\delta_{\ov{s}}(f-S_{n}^{(\gamma)}(f))\|_{X}=0,\notag \\


%


\sum_{\la\ov{s},\ov{\gamma}\ra>\nu} 


\|\delta_{\ov{s}}(f-S_{n}^{(\ov{\gamma})}(f))\|_{X}= 


\sum_{\la\ov{s},\ov{\gamma}\ra\ge n}\|\delta_{\ov{s}}(f)\|_{X}


\end{gather}


для всех $\nu=1,\dotsc,n$.





Пусть $\nu> n$, тогда 


\begin{gather}


\sum_{\la\ov{s},\ov{\gamma}\ra\le\nu}\frac{1}{\varphi(p^{-\|\ov{s}\|})} 


\|\delta_{\ov 


s}(f-S_{n}^{\ov{(\gamma)}}(f))\|_{X}=\sum_{n<\la\ov{s},\ov{\gamma}\ra\leq\nu} 


\frac{1}{\varphi(p^{-\|\ov{s}\|})}\|\delta_{\ov{s}}(f)\|_{X},\\


%


\sum_{\la\ov{s},\ov{\gamma}\ra>\nu}\|\delta_{\ov{s}}(f-S_{n}^{(\ov{\gamma})}) 


\|_{X}=\sum_{\la\ov{s},\ov{\gamma}\ra>\nu}\|\delta_{\ov{s}}(f)\|_{X}.


\end{gather}


Из соотношений (14)--(17) следует


\begin{multline}


\sigma_{\nu}(f)=\frac{p^{-\nu/q}}{\varphi(p^{-\nu})} 


\sum_{\la\ov{s},\ov{\gamma}\ra>\nu}\|\delta_{\ov{s}}(f-S_{n}^{(\ov{\gamma})}) 


\|_{X}+p^{-\nu/q}\sum_{0<\la\ov{s},\ov{\gamma}\ra\le\nu} 


\frac{1}{\varphi(p^{-\|\ov{s}}\|)} 


\|\delta_{\ov{s}}(f-S_{n}^{(\ov{\gamma})})\|_{X}\le \\


%


\le C


\begin{cases}


\frac{p^{-\nu/q}}{\varphi(p^{-\nu})}\sum\limits_{\la\ov{s},\ov{\gamma}\ra 


>n}\|\delta_{\ov{s}}(f)\|_{X} & \text{при} \ \ \nu\le n,\\[2mm]


\frac{p^{-\nu/q}}{\varphi(p^{-\nu})}\sum\limits_{\la\ov{s},\ov{\gamma}\ra 


>\nu}\|\delta_{\ov{s}}(f)\|_{X}+p^{-\nu/q}\sum\limits_{n< 


\la\ov{s},\ov{\gamma}\ra\le\nu}\frac{1}{\varphi(p^{-\|\ov{s}\|})} 


\|\delta_{\ov{s}}(f)\|_{X} & \text{при} \ \ \nu>n.


\end{cases}


\end{multline}  


По условию теоремы $0<r_{1}<\log_{2}\alpha_{\varphi}$. Поэтому существует 


$\Phi$-функция $\theta_{1}(t)$ такая, что


$\theta_{1}(t)\asymp\frac{\varphi(t)}{t^{r_{1}+\varepsilon}}$, 


$\varepsilon>0$ (см. лемму~А). Следовательно, учитывая, что 


$\theta_{1}(t)\uparrow$ и пользуясь леммой В, получим 


\begin{multline}


\sum_{n<\la\ov{s},\ov{\gamma}\ra\le\nu}


\frac{1}{\varphi(p^{-\|\ov{s}\|})}\|\delta_{\ov{s}}(f)\|_{X}


\le C\sum_{n<\la\ov{s},\ov{\gamma}\ra\le\nu}\frac{1}{\varphi(p^{-\|\ov{s}\|


})}p^{-\la\ov{s},\ov{r}\ra}= \\


%


=C\sum_{n< \la\ov{s},\ov{\gamma}\ra\le\nu}\frac{1}{\varphi(p^{-\|\ov{s}\|})} 


p^{-r_{1}\la\ov{s},\ov{\gamma}\ra}\le C\sum_{n< 


\la\ov{s},\ov{\gamma}\ra\le\nu}\frac{1}{\varphi(p^{-\la\ov{s},\ov{\gamma}\ra})} 


p^{-r_{1}\la\ov{s},\ov{\gamma}\ra}\le \\


%


\le C\sum_{n< \la\ov{s},\ov{\gamma}\ra\le\nu}\theta_{1}^{-1}(p^{-\la\ov{s},\ov


{\gamma}\ra})\le 


C\frac{1}{\theta_{1}(p^{\nu})}\sum_{n< \la\ov{s},\ov{\gamma}\ra\le\nu}p^{


\varepsilon\la\ov{s},\ov{\gamma}\ra}\le C\frac{1}{\theta_{1}(p^{-\nu})}p^{\nu


\varepsilon}\nu^{d-1}\le


C\frac{p^{-\nu r_{1}}}{\varphi(p^{-\nu})}\nu^{d-1}


\end{multline}


для любой функции $f$ из $H_{X}^{\ov{r}}$. Далее, по лемме Б имеем


\begin{equation}


\sum_{\la\ov{s},\ov{\gamma}\ra >\nu}\|\delta_{\ov{s}}(f)\|_{X}\le C\sum_{\la\ov


{s},\ov{\gamma}\ra >\nu}p^{-\la\ov{s},\ov{r}\ra}\le Cp^{-\nu 


r_{1}}\nu^{d-1}.


\end{equation}  


Из неравенств (18)--(20) для любой функций $f\in H_{X}^{\ov{r}}$ получим


\begin{equation}


\sigma_{\nu}(f)\le C


\begin{cases}


\frac{p^{-\nu/q}}{\varphi(p^{-\nu})}p^{-nr_{1}}\nu^{d-1}


& \text{при} \ \ \nu\le n, \\


\frac{p^{-\nu(r_1+\frac{1}{q})}}{\varphi(p^{-\nu})}\nu^{d-1}


& \text{при} \ \ \nu>n.


\end{cases}


\end{equation}  








По условию теоремы $\frac{1}{q}<\log_{2}\varphi$. Следовательно, по лемме


А существует $\Phi$-функция $\theta_{1}(t)$ такая, что 


$\frac{\varphi(t)}{t^{\frac{1}{q}}}\asymp\theta_{1}(t)\uparrow$. Поэтому 


\begin{equation}


\frac{p^{-\nu/q}}{\varphi(p^{-\nu})}\le C\frac{p^{-n/q}}{\varphi(p^{-n})}, 


\quad \nu\le n.  


\end{equation}   


Так как $\beta_{\varphi}<2^{\frac{1}{q}+r_{1}}$ по условию теоремы, то 


существует число $\varepsilon\in (0,r_{1}+\frac{1}{q})$  такое, что 


$\beta_{\varphi}<2^{\frac{1}{q}+r_{1}-\varepsilon}<2^{\frac{1}{q}+r_{1}}$. 


Полагая $\gamma(t)=t^{r_{1}+\frac{1}{q}-\varepsilon}$ в лемме A,  


легко убеждаемся в существовании $\Phi$-функ\-ции $\theta_{1}(t)$ такой, что 


$\frac{t^{\frac{1}{q}+r_{1}-\varepsilon}}{\varphi(t)}\asymp 


\theta_{1}(t)\uparrow$. Следовательно,


\begin{equation}


\frac{p^{\nu(r_{1}+\frac{1}{q})}\nu^{d-1}}{\varphi(p^{-\nu})}\le


C\theta_{1}(p^{-\nu})\nu^{d-1}\le C\theta_{1}(p^{-n})


\frac{n^{d-1}}{p^{n\varepsilon}}\le


C\frac{p^{-n(\frac{1}{q}+r_{1})}}{\varphi(p^{-n})}n^{d-1}


\end{equation}  


для $\nu>n$. Из соотношений (21)--(23) следует


$$


\sup_{\nu>0}\sigma_{\nu}(f)\le C\frac{p^{-n(\frac{1}{q}+r_{1})}}{\varphi


(p^{-n})}n^{d-1}


$$ 


для любой из функций $f\in H_{X}^{\ov{r}}$. Таким образом (см. (14)),


$$


\sup_{f\in H_{X}^{\ov{r}}}\|f-S_{n}^{(\ov{\gamma})}(f)\|_{q,\infty}\le 


C(d,p)\frac{p^{-n(\frac{1}{q}+r_{1})}}{\varphi(p^{-n})}n^{d-1}.


$$ 








\begin{thm}


Пусть обобщенные системы Хаара $\chi\{p_{n}^{(j)}\}$, 


$j=1,\dotsc,d$, определены одним натуральным числом


$p_{n}^{(j)}=p\geq 2$, $n=1,2,\dotsc;$ $j=1,\dotsc,d$, и  выполняются 


неравенства $\max\big\{\frac{1}{q},r\big\}<\frac{1}{q_{1}}<\frac{1}{q}+r$. 


Тогда имеет место соотношение


$$


E_{Q_{n}}(H_{L_{q_{1}}}^{r})_{q,\infty}\equiv \sup_{f\in 


H_{L_{q_{1}}}^{r}}\|f-S_{n}(f)\|_{q,\infty}\asymp 


p^{-n(r+\frac{1}{q}-\frac{1}{q_{1}})}\cdot n^{d-1}.


$$


\end{thm}








\begin{pf}


Оценка сверху величины


$E_{Q_{n}}(H_{L_{q_{1}}}^{r})_{q,\infty}$ следует из теоремы 1.





Для оценки снизу рассмотрим произведение $f_{0}(\ov x)= 


\prod\limits_{j=1}^{d}f_{j}^{0}(x_{j})$ функций


$f^0_{j}(x_{j})$, построенных в 


лемме 4, в случае пространства Лебега $X(\varphi)=L_{q_{1}}[0,1]$, т.\,е. 


$\varphi(t)=t^{\frac{1}{q_{1}}}$. Согласно этой лемме $f_{0}\in 


H_{L_{q_{1}}}^{r}$. По определению точной верхней грани для всех $\ov 


l=(l_{1},\dotsc,l_{d})$ таких, что $\|\ov l\|=n$, имеет место неравенство


\begin{multline}


\|f_{0}-S_{n}(f_{0})\|_{q,\infty}\geq p^{(1-\frac{1}{q})\|\ov l\|} 


\int_{0}^{p^{-l_{1}}}\cdots\int_{0}^{p^{-l_{d}}}|f_{0}(\ov 


x)-S_{n}(f_{0},\ov x)|d\ov x= \\


%


=p^{n(1-\frac{1}{q})} \sum_{\nu=n}^{\infty}\;\sum_{\nu<\|\ov s\|\leq\nu+1}


\int_{p^{-s_{1}-1}}^{p^{-s_{1}}}\cdots\int_{p^{-s_{d}-1}}^{p^{-s_{d}}}|f_{0} 


(\ov x)-S_{n}(f_{0},\ov x)|d\ov x.


\end{multline}


Так как носители функций $\chi_{p^{k_{j}}+p}(x_{j})$, $j=1,\dotsc,d$, не 


пересекаются, то 


$$


\int_{p^{-s_{1}-1}}^{p^{-s_{1}}}\dotsb\int_{p^{-s_{d}-1}}^{p^{-s_{d}}}|f_{0} 


(\ov x)-S_{n}(f_{0},\ov x)|d\ov x=p^{-(r+1)\|\ov 


s\|}\cdot p^{\frac{\|\ov s\|}{q_{1}}}.


$$


Поэтому из неравенства (24) следует


\begin{equation}


\|f_{0}-S_{n}(f_{0})\|_{q,\infty}\geq p^{n(1-\frac{1}{q})} 


\sum_{\nu=n}^{\infty}\;\sum_{\nu<\|\ov s\|\leq\nu+1}


p^{-(r+1)\|\ov s\|}\cdot p^{\frac{\|\ov s\|}{q_{1}}}.


\end{equation}


Пользуясь леммой Б, получим 


\begin{equation}


\sum_{\nu=n}^{\infty}\sum_{\nu<\|\ov s\|\leq\nu+1} 


p^{-(r+1-\frac{1}{q_{1}})\|\ov s\|}\geq 


C(r)\cdot p^{-n(r+1-\frac{1}{q_{1}})}\cdot n^{d-1}.


\end{equation}


Из неравенств (25), (26) вытекает, что


$\sup\limits_{f\in H_{L_{q_{1}}}^{r}}\|f-S_{n}(f)\|_{q,\infty}\geq 


p^{-n(r+\frac{1}{q}-\frac{1}{q_{1}})}\cdot n^{d-1}$.


\end{pf}








\begin{thebibliography}{99}





\bibitem{1}


Крейн С.Г., Петунин Ю.И., Семенов Е.М. {\em Интерполяция 


линейных операторов}. -- М.: Наука, 1978. -- 400~с.


\bibitem{2}


Виленкин Н.Я. {\em Об одном классе полных ортонормированных 


систем} // Изв. АН СССР. Сер. матем. -- 1947. -- T.\,11. -- С.\,363--400.


\bibitem{3}


Голубов Б.И. {\em Об одном классе полных ортогональных систем} 


// Сиб. матем. журн. -- 1968. -- T.\,9. -- С.\,297--314.


\bibitem{4}


Андрианов А.В. {\em Приближение функций из классов 


$MH_{q}^{r}$ полиномами Хаара} // Матем. заметки. --


1999. -- T.\,66. -- С.\,323--335.


\bibitem{5}


Кашин Б.С. {\em О нижних оценках для $n$-членных приближений} 


// Матем. заметки. -- 2001. -- T.\,70. -- С.\,636--638. 


\bibitem{6}


Temlyakov V.N. {\em Nonlinear $m$-term approximation with 


regard to the multivariate Haar system} // East J. Approx. --


1998. -- V.\,4. -- P.\,87--106.


\bibitem{7}


Освальд П. {\em Об $N$-членных приближениях по системе Хаара 


в $H^{s}$-нормах} // Метрич. теория функ. и смежные вопр. анал. -- М., 1999. 


-- С.\,137--163.


\bibitem{8}


Лапин С.В. {\em Некоторые теоремы вложения для произведений 


функций} // Деп. в ВИНИТИ, 1980, \No\,1036-80. -- 31~с.


\bibitem{9}


Темляков В.Н. {\em Приближение функций с ограниченной 


смешанной производной} // Тр. МИАН. -- 1986. -- T.\,178. -- 112~с.


\bibitem{10}


Борисова Е.А. {\em О вложении классов функций, заданных 


последовательностями наилучших приближений по системам типа Хаара} // 


Деп. в ВИНИТИ, 1986, \No\,6333-86. -- 52~с.


\bibitem{11}


Акишев Г.А., Коспанова Р.К. {\em Теоремы вложения в 


симметричные пространства и обобщенная система Хаара} // Деп.


в КазНИИНТИ, 1991, \No\,3618. -- 50~с.





\end{thebibliography}





\vskip 2cm





\postup{\it Карагандинский государственный}{\qquad \it Поступили}


\postup{\it университет}{\it первый вариант $28.10.2002$}


\postup{}{\it окончательный вариант $11.08.2003$}





\label{end}


\end{document}


